\section{Cadre théorique et épistémologique}\label{sec:cadre}

Pour analyser les tensions introduites par les IAG, il convient de préciser les présupposés épistémologiques sur lesquels repose l'évaluation scolaire.
Ces présupposés ne sont pas des détails techniques: ils constituent le fondement implicite de ce que le système scolaire entend par \emph{savoir}, \emph{compétence} et \emph{réussite}.
C'est précisément parce qu'ils sont rarement explicités qu'ils se révèlent fragiles dès lors qu'un outil vient les mettre à l'épreuve.

\subsection{La transposition didactique et ses effets épistémologiques}

La théorie de la \emph{transposition didactique}, développée par Chevallard (1985), part d'un constat simple mais aux conséquences profondes: le savoir qui circule à l'école n'est pas le savoir produit par les savants. 
Il en est une version transformée, décontextualisée puis recontextualisée, rendue enseignable au prix d'un éloignement de ses conditions réelles de production. 
Une démonstration mathématique telle qu'elle figure dans un manuel n'est pas la démonstration telle qu'elle a été construite par un chercheur: elle a été épurée de ses tâtonnements, de ses erreurs, de ses dépendances à un contexte particulier.

Ce détachement a une conséquence directe sur l'évaluation: on évalue non pas une capacité à produire du savoir dans un contexte réel, mais une maîtrise d'une version scolaire, normée et stabilisée de ce savoir.
L'approche par compétences cherche précisément à corriger cet écart, en réintroduisant des situations-problèmes représentatives de la vie réelle. 
Mais elle se heurte alors à une difficulté nouvelle: plus les tâches évaluées ressemblent à des situations réelles, plus elles sont exposées à l'intervention d'outils réels --- comme les IAG.\@
La transposition didactique avait jusqu'ici protégé l'évaluation scolaire en la maintenant dans un espace artificiel; en voulant sortir de cet espace, l'approche par compétences expose ses propres présupposés.

\subsection{La reproduction sociale et la construction de la légitimité scolaire}

Le cadre posé par Bourdieu apporte un second niveau d'analyse, distinct mais
complémentaire. 
Pour Bourdieu et Passeron (\citeyear{pluet1971pierre}), le système scolaire n'est pas un dispositif neutre de transmission des savoirs: il est un mécanisme de reproduction sociale qui transforme des inégalités d'origine en différences de mérite. 
L'école consacre comme légitime un certain rapport au savoir --- distancié, formel, décontextualisé --- qui est précisément celui que les familles cultivées transmettent à leurs enfants. 
Ce faisant, elle masque le caractère arbitraire de cette légitimité derrière l'apparence de l'objectivité des critères.

Ce mécanisme a des implications directes pour notre sujet. 
Si les critères d'évaluation scolaire valorisent un style d'écrit, une structure argumentative, une maîtrise lexicale que certains héritent et que d'autres doivent laborieusement construire, alors ces critères ne mesurent pas seulement une compétence: ils mesurent aussi une familiarité culturelle.
Les IAG perturbent ce dispositif de manière ambivalente. 
D'un côté, elles pourraient réduire l'inégalité en donnant accès à des productions formellement conformes aux attentes institutionnelles à des apprenants qui n'ont pas reçu ce capital culturel en héritage. 
De l'autre, elles creusent une nouvelle inégalité entre ceux qui savent exploiter ces outils de façon critique et ceux qui s'y soumettent sans discernement. 
Dans les deux cas, elles mettent en évidence ce qu'on pourrait appeler l'arbitraire des critères: si une machine peut les satisfaire, c'est qu'ils ne mesurent peut-être pas ce qu'ils prétendent mesurer.

\subsection{L'épistémologie de l'évaluation: présupposés et fragilités}

L'évaluation scolaire repose sur un ensemble de présupposés qui n'ont jamais eu besoin d'être formulés explicitement parce qu'ils semblaient aller de soi.%
~\cite{leroux2015evaluer} en identifie plusieurs dans son analyse des pratiques d'évaluation au collégial: que le travail remis est le produit authentique de l'apprenant; qu'il est représentatif de ses capacités réelles; que la notation est cohérente et comparable entre contextes différents. 
Ces présupposés étaient déjà fragiles avant les IAG: la variabilité inter-correcteurs est bien documentée, les biais sociaux dans l'évaluation sont établis. % chktex 13
Mais ils restaient des fragilités contenues, car la production individuelle écrite demeurait difficile à falsifier sans laisser de traces.

Les IAG effondrent cette dernière barrière. 
Elles peuvent générer un travail formellement conforme aux critères d'évaluation sans que l'apprenant ait engagé les processus cognitifs que l'évaluation cherchait à stimuler.
La notion centrale de \emph{validité de l'évaluation} --- sa capacité à mesurer ce qu'elle prétend mesurer --- est ainsi mise en crise: non pas fragilisée à la marge, mais questionnée dans son principe. 
Le cadre de Chevallard permet de comprendre pourquoi cette crise est structurelle et non accidentelle; le cadre de Bourdieu permet de comprendre pourquoi
elle est aussi une crise de légitimité, pas seulement d'efficacité technique.